For every fixed residual tuple and \(0<c<1\), there are constants \(0<C_{\mathrm{low}}\le C_{\mathrm{up}}<\infty\) and \(n_0\in\mathbb N\), independent of \(n\), such that \[C_{\mathrm{low}}n^{-r_{\mathrm{low}}}\le\mathcal M_n(\alpha,\beta,s,d,M_Y,L,c,t_0)\le C_{\mathrm{up}}n^{-R_{\mathrm{ach}}}\qquad(n\ge n_0),\] where \[r_{\mathrm{low}}=\min\left\{\frac{2\alpha}{2\alpha+1},\frac{8}{4+2/(\alpha\vee\beta)+d/s}\right\}.\] Hence \(R_{\mathrm{ach}}\le r_{\mathrm{low}}\), and any exact polynomial exponent must satisfy \[R_{\mathrm{ach}}\le\rho(\alpha,\beta,s,d)\le r_{\mathrm{low}}.\] At every tuple for which \(R_{\mathrm{ach}}=r_{\mathrm{low}}\), the exact exponent exists and equals that common value; it is \(\beta\)-invariant on any subregion where the equality and common value persist as \(\beta\) varies. Separately, at the excluded boundary \(c=1\), Lemma~\texttt{lem:unit-positivity-oracle-rate} proves the exact exponent \(2\alpha/(2\alpha+1)\) for every \(\beta,s,d\).